\documentclass[12pt]{article}
\usepackage[margin=0.8in]{geometry}
\usepackage{amsmath}
\usepackage{xcolor}
\usepackage{tcolorbox}
\usepackage{enumitem}
\usepackage{tikz}

\title{\textbf{Convex Preferences: What Q3b Is Really Asking}}
\author{}
\date{}

\begin{document}
\maketitle

\section*{The Setup (from Q3b)}

Melsi has preferences: $u(C_1, C_2) = v(C_1) + v(C_2)$

where $v$ is \textbf{increasing and concave} (like $v(x) = \ln(x)$ or $v(x) = \sqrt{x}$).

The question asks: "Are her preferences over bundles $(C_1, C_2)$ convex?"

\section*{What Does "Convex Preferences" Mean?}

\begin{tcolorbox}[colback=blue!5!white,colframe=blue!75!black]
\textbf{Simple Definition:}

Convex preferences mean you prefer \textbf{averages to extremes}.

If you like bundle A and you like bundle B, then you like the \textbf{mixture} (half of A, half of B) at least as much as A or B alone.

\vspace{0.3cm}

\textbf{In pictures:} Your indifference curves are "bowed toward the origin" (convex to the origin).
\end{tcolorbox}

\subsection*{IMPORTANT: Convex vs Concave Preferences}

\begin{tcolorbox}[colback=red!10!white,colframe=red!75!black,title=Don't Get These Confused!]
\textbf{Convex Preferences} (what we usually assume):
\begin{itemize}
\item You prefer \textbf{balanced} bundles (averages)
\item Indifference curves are \textbf{bowed toward the origin}
\item You have \textbf{diminishing} marginal rate of substitution
\item As you get more of one good, you're willing to give up \textit{less} of the other good to get even more
\end{itemize}

\vspace{0.3cm}

\textbf{Concave Preferences} (rare and weird):
\begin{itemize}
\item You prefer \textbf{extreme} bundles (all of one thing)
\item Indifference curves are \textbf{bowed away from the origin}
\item You have \textbf{increasing} marginal rate of substitution
\item The more you have of one good, the MORE you're willing to give up of the other to get even more
\item Example: "The more coffee I drink, the MORE I'm willing to sacrifice donuts for additional coffee" (addiction-like behavior)
\end{itemize}

\vspace{0.3cm}

\textbf{Key takeaway:} Most goods have \textbf{convex preferences} because people like variety and balance. Concave preferences don't make economic sense for typical goods.
\end{tcolorbox}

\subsection*{Intuitive Example}

Suppose you have two bundles:
\begin{itemize}
\item Bundle A: 10 cookies today, 0 cookies tomorrow
\item Bundle B: 0 cookies today, 10 cookies tomorrow
\end{itemize}

With convex preferences, you'd prefer:
\begin{itemize}
\item Bundle C: 5 cookies today, 5 cookies tomorrow (the average)
\end{itemize}

\textbf{Why?} Because you like balance! Having all your cookies on one day is extreme. Spreading them out is better.

\subsection*{The Graphical Picture}

\begin{center}
\begin{tikzpicture}[scale=1.2]
% Axes
\draw[->] (0,0) -- (5,0) node[right] {$C_1$ (today)};
\draw[->] (0,0) -- (0,5) node[above] {$C_2$ (tomorrow)};

% Indifference curve (convex)
\draw[thick, blue] (0.5,4.5) .. controls (1.5,2) and (2,1.5) .. (4.5,0.5);
\node[blue] at (4,2.5) {Indifference curve};

% Points
\filldraw[red] (1,3) circle (2pt) node[above left] {A};
\filldraw[red] (3,1) circle (2pt) node[below right] {B};
\filldraw[green!50!black] (2,2) circle (2pt) node[above right] {Average (better!)};

% Line between A and B
\draw[dashed, red] (1,3) -- (3,1);

% Arrow showing preference
\draw[->, very thick, green!50!black] (2.3,2.3) -- (2.5,2.5);

\end{tikzpicture}
\end{center}

The average bundle is on a \textbf{higher} indifference curve than A or B!

\newpage

\section*{What is MU$_1$/MU$_2$?}

\begin{tcolorbox}[colback=yellow!10!white,colframe=orange!75!black]
\textbf{MU = Marginal Utility}

\begin{itemize}
\item $MU_1 = \frac{\partial u}{\partial C_1}$ = how much extra utility you get from one more unit of good 1
\item $MU_2 = \frac{\partial u}{\partial C_2}$ = how much extra utility you get from one more unit of good 2
\end{itemize}

\vspace{0.3cm}

\textbf{The ratio MU$_1$/MU$_2$:}

This is your \textbf{marginal rate of substitution (MRS)}!

\[
MRS = -\frac{dC_2}{dC_1}\Big|_{IC} = \frac{MU_1}{MU_2}
\]

It tells you: "How many units of good 2 would I give up for one more unit of good 1?"
\end{tcolorbox}

\subsection*{In Q3b}

For $u(C_1, C_2) = v(C_1) + v(C_2)$:

\begin{align*}
MU_1 &= \frac{\partial u}{\partial C_1} = v'(C_1) \\
MU_2 &= \frac{\partial u}{\partial C_2} = v'(C_2)
\end{align*}

So:
\[
MRS = \frac{MU_1}{MU_2} = \frac{v'(C_1)}{v'(C_2)}
\]

\section*{The Convexity Test}

\begin{tcolorbox}[colback=green!5!white,colframe=green!75!black]
\textbf{Preferences are convex if and only if:}

As you move along an indifference curve from left to right (increasing $C_1$, decreasing $C_2$), the MRS \textbf{decreases} (becomes less steep).

\vspace{0.3cm}

In other words: $\frac{MU_1}{MU_2}$ should get smaller as $C_1$ increases and $C_2$ decreases.
\end{tcolorbox}

\subsection*{Why This Makes Sense}

\textbf{Start:} You have lots of $C_2$, very little $C_1$

$\Rightarrow$ You really want more $C_1$ (it's scarce!)

$\Rightarrow$ MRS is high (willing to give up lots of $C_2$ for one unit of $C_1$)

\vspace{0.3cm}

\textbf{End:} You have lots of $C_1$, very little $C_2$

$\Rightarrow$ You already have plenty of $C_1$, now you want $C_2$

$\Rightarrow$ MRS is low (not willing to give up much $C_2$ for more $C_1$)

\vspace{0.3cm}

This is \textbf{diminishing marginal rate of substitution} = convex preferences!

\newpage

\section*{Solving Q3b Step-by-Step}

\textbf{Given:} $u(C_1, C_2) = v(C_1) + v(C_2)$ where $v$ is \textbf{concave}

\textbf{Question:} Are preferences convex?

\subsection*{Step 1: Find MRS}

\[
MRS = \frac{MU_1}{MU_2} = \frac{v'(C_1)}{v'(C_2)}
\]

\subsection*{Step 2: What happens when $C_1 \uparrow$ and $C_2 \downarrow$?}

As we move along an IC:
\begin{itemize}
\item $C_1$ increases: $C_1 \uparrow$
\item $C_2$ decreases: $C_2 \downarrow$
\end{itemize}

\subsection*{Step 3: Use the fact that $v$ is concave}

\textbf{Concave function means:} $v'' < 0$

This means $v'$ is \textbf{decreasing}:
\begin{itemize}
\item As $C_1 \uparrow$, we have $v'(C_1) \downarrow$ (marginal utility of $C_1$ falls)
\item As $C_2 \downarrow$, we have $v'(C_2) \uparrow$ (marginal utility of $C_2$ rises)
\end{itemize}

\subsection*{Step 4: What happens to MRS?}

\[
MRS = \frac{v'(C_1)}{v'(C_2)}
\]

\begin{itemize}
\item Numerator $v'(C_1)$ goes down: $v'(C_1) \downarrow$
\item Denominator $v'(C_2)$ goes up: $v'(C_2) \uparrow$
\end{itemize}

\begin{tcolorbox}[colback=yellow!20!white,colframe=orange!75!black]
\textbf{Result:}

\[
\frac{v'(C_1) \downarrow}{v'(C_2) \uparrow} \quad \Rightarrow \quad MRS \downarrow
\]

The MRS decreases! This is exactly what we need for convex preferences.

\vspace{0.3cm}

\textbf{Answer: YES, preferences are convex.}
\end{tcolorbox}

\newpage

\section*{Concrete Example}

Let's use $v(x) = \ln(x)$ (which is concave).

So $v'(x) = \frac{1}{x}$ (which is decreasing).

\subsection*{Scenario 1: Lots of $C_2$, little $C_1$}

Suppose $C_1 = 1$, $C_2 = 10$:

\begin{align*}
v'(C_1) &= v'(1) = \frac{1}{1} = 1 \\
v'(C_2) &= v'(10) = \frac{1}{10} = 0.1 \\
MRS &= \frac{1}{0.1} = 10
\end{align*}

\textbf{Interpretation:} You'd give up 10 units of $C_2$ for 1 more unit of $C_1$ (because $C_1$ is scarce!).

\subsection*{Scenario 2: Lots of $C_1$, little $C_2$}

Suppose $C_1 = 10$, $C_2 = 1$:

\begin{align*}
v'(C_1) &= v'(10) = \frac{1}{10} = 0.1 \\
v'(C_2) &= v'(1) = \frac{1}{1} = 1 \\
MRS &= \frac{0.1}{1} = 0.1
\end{align*}

\textbf{Interpretation:} Now you'd only give up 0.1 units of $C_2$ for 1 more unit of $C_1$ (because $C_2$ is now scarce!).

\subsection*{The Pattern}

\[
MRS \text{ went from } 10 \text{ down to } 0.1 \quad \Rightarrow \quad \text{MRS is decreasing} \quad \Rightarrow \quad \text{Convex!}
\]

\section*{Why Does This Matter?}

\begin{enumerate}[leftmargin=*]
\item \textbf{Convexity ensures interior solutions exist}
\begin{itemize}
\item With convex preferences, you want to balance your consumption
\item You won't go to extremes (all $C_1$ or all $C_2$)
\end{itemize}

\item \textbf{Convexity is realistic}
\begin{itemize}
\item Most people prefer variety/balance
\item "Diminishing marginal utility" is a real phenomenon
\end{itemize}

\item \textbf{Mathematically nice}
\begin{itemize}
\item Guarantees unique optimal solution
\item FOC is sufficient (not just necessary)
\end{itemize}
\end{enumerate}

\newpage

\section*{Summary}

\begin{tcolorbox}[colback=green!5!white,colframe=green!75!black]
\textbf{Key Concepts:}

\begin{enumerate}[leftmargin=*]
\item \textbf{Convex preferences:} Prefer averages to extremes; ICs are bowed toward origin

\item \textbf{MRS = MU$_1$/MU$_2$:} The slope of your indifference curve; how willing you are to trade good 2 for good 1

\item \textbf{Convexity test:} MRS should decrease as you move along IC (as $C_1 \uparrow$, $C_2 \downarrow$)

\item \textbf{For Q3b:}
\begin{itemize}
\item $v$ concave $\Rightarrow$ $v'$ decreasing
\item As $C_1 \uparrow$: $v'(C_1) \downarrow$ (numerator falls)
\item As $C_2 \downarrow$: $v'(C_2) \uparrow$ (denominator rises)
\item Therefore: $\frac{v'(C_1)}{v'(C_2)} \downarrow$ (MRS falls)
\item Conclusion: Preferences are convex!
\end{itemize}
\end{enumerate}
\end{tcolorbox}

\section*{The Connection}

\begin{center}
\fbox{\begin{minipage}{0.9\textwidth}
\textbf{Concave $v$ function}

$\Downarrow$

Marginal utility $v'$ is decreasing

$\Downarrow$

MRS = $\frac{v'(C_1)}{v'(C_2)}$ decreases along IC

$\Downarrow$

\textbf{Convex preferences}

$\Downarrow$

You prefer balanced bundles to extreme bundles
\end{minipage}}
\end{center}

\section*{Common Mistakes to Avoid}

\begin{itemize}[leftmargin=*]
\item \textbf{Don't confuse:} "concave function" vs "convex preferences"
\begin{itemize}
\item In Q3b: $v$ is \textbf{concave}, but preferences are \textbf{convex}!
\item Concave $v$ $\Rightarrow$ diminishing marginal utility $\Rightarrow$ convex preferences
\end{itemize}

\item \textbf{Remember the direction:} When checking convexity, you need MRS to \textbf{decrease} as $C_1$ increases

\item \textbf{Be clear about notation:} MRS can be written as:
\begin{itemize}
\item $MRS = -\frac{dC_2}{dC_1}$ (negative slope of IC)
\item $|MRS| = \frac{MU_1}{MU_2}$ (absolute value)
\end{itemize}
Sometimes people drop the negative sign for simplicity.
\end{itemize}

\end{document}
